\documentclass[11pt]{article}
\usepackage[T1]{fontenc}
\usepackage{lmodern}
\usepackage[margin=1.04in]{geometry}
\usepackage{amsmath,amssymb,amsthm,mathrsfs,booktabs,microtype}
\usepackage[colorlinks=true,linkcolor=blue,citecolor=blue,urlcolor=blue]{hyperref}
\numberwithin{equation}{section}
\newtheorem{theorem}{Theorem}[section]
\newtheorem{proposition}[theorem]{Proposition}
\newtheorem{lemma}[theorem]{Lemma}
\newtheorem{corollary}[theorem]{Corollary}
\theoremstyle{remark}\newtheorem{remark}[theorem]{Remark}
\newcommand{\Q}{\mathbb Q}\newcommand{\Z}{\mathbb Z}
\newcommand{\Qbar}{\overline{\mathbb Q}}
\newcommand{\eps}{\varepsilon}\newcommand{\LL}{\mathcal L}
\newcommand{\HH}{\mathcal H}\newcommand{\wt}{\widehat}
\DeclareMathOperator{\spanop}{span}
\title{A half-exponent bound for the block complexity\\of algebraic numbers}
\author{}\date{}
\hypersetup{pdftitle={A half-exponent bound for the block complexity of algebraic numbers},pdfsubject={Research manuscript: moving weights and expansions in integer bases}}
\begin{document}\maketitle
\begin{abstract}
Let $\xi$ be a real algebraic irrational number and $b\ge2$ an integer.
Let $p(L)=p_{\xi,b}(L)$ count the distinct blocks of length $L$ in the
fractional base-$b$ expansion of $\xi$. We prove
\[
 \limsup_{L\to\infty}
 \frac{p(L)(\log\log L)^\gamma}{L\sqrt{\log L}}=+\infty
 \qquad(\gamma>1).
\]
Consequently $p(L)/(L(\log L)^\eta)$ is unbounded for every $\eta<1/2$.
The new input is a sparse-parameter bound for semistable twisted heights
in dimension three with fixed forms and fixed finite support, while the
weights may vary from one parameter to the next. We prove this bound by
adapting the auxiliary-polynomial argument of Evertse and Ferretti.
The number of strongly separated parameters is $O(\delta^{-2})$ and
the uniform cutoff has double logarithm
$O(\delta^{-2}\log(2/\delta))$. An explicit calculation verifies
semistability for the full family of weights arising from repetitions.
The argument includes nonreal conjugates and composite bases.
\end{abstract}

\section{Introduction}
Fix a real algebraic irrational number $\xi$ and an integer $b\ge2$.
Replacing $\xi$ by its fractional part leaves its block complexity
unchanged, so we henceforth assume $0<\xi<1$.
Write $\xi=\sum_{i\ge1}d_ib^{-i}$, where $d_i\in\{0,\ldots,b-1\}$, and put
\[
 p(L)=p_{\xi,b}(L)=\#\{d_{j+1}\cdots d_{j+L}:j\ge0\}.
\]
We use natural logarithms, and expressions involving iterated logarithms
are restricted to sufficiently large arguments.

\begin{theorem}\label{thm:main}
For every real algebraic irrational $\xi$, every integer $b\ge2$,
and every real $\gamma>1$,
\begin{equation}\label{eq:main}
 \limsup_{L\to\infty}
 \frac{p(L)(\log\log L)^\gamma}{L(\log L)^{1/2}}=+\infty.
\end{equation}
In particular, for every real $\eta<1/2$,
\begin{equation}\label{eq:eta}
 \limsup_{L\to\infty}\frac{p(L)}{L(\log L)^\eta}=+\infty.
\end{equation}
\end{theorem}

These are limsup statements. Neither the unweighted endpoint
$\eta=1/2$ nor the case $\gamma=1$ is asserted.
Adamczewski and Bugeaud proved $p_{\alpha,b}(L)/L\to\infty$ for algebraic
irrational $\alpha$ and integer bases $b\ge2$ \cite{AB07}.
Bugeaud and Evertse obtained the logarithmic limsup refinement with
$\eta<1/11$ \cite[Theorem~2.1]{BE08}. Our argument uses their
repetition method and the semistable machinery in \cite[Sections~8--14]{EF13}.

The essential issue is the ratio $\rho=r/t$ in an approximation
\[
 |(b^t-b^r)\xi-a|\le b^{-v}.
\]
A partition of $\rho$ into intervals of length comparable to $\eps$
permits fixed weights, but costs a factor of order $\eps^{-1}$.
We instead retain the exact value of $\rho$ at each point. This requires
a separate argument: the fixed-weight statement of
\cite[Theorem~8.1]{EF13} does not itself allow this substitution.
Section~\ref{sec:moving} proves the needed variant, including its cutoff.
Its use here depends on an explicit semistability calculation in
Section~\ref{sec:stable}.

All constants in the moving-weight argument depend only on the fixed
number field, forms, and finite set of places. They are independent of
the individual weights and parameters. No effective numerical starting
length for \eqref{eq:main} is claimed. The exponent range is the same
for every fixed pair $(\xi,b)$; the constants and cutoffs may depend
on that pair.

\section{Repetitions and a sparse sequence}\label{sec:words}
Fix $u>0$, $\gamma\ge0$, and set
\[
 F(x)=(\log x)^u(\log\log x)^{-\gamma}.
\]
This function eventually increases to infinity, satisfies $F(x)=o(x)$,
and, for every fixed $D>0$,
\begin{equation}\label{eq:slow}
 F(\lceil DxF(x)\rceil)/F(x)\longrightarrow1.
\end{equation}
Indeed, its logarithmic derivative is
$(u-\gamma/\log\log x)/(x\log x)$; also
$\log\lceil DxF(x)\rceil=\log x+O(\log\log x)$.
Suppose, for a contradiction to be specified at the end, that
\begin{equation}\label{eq:LC}
 p(k)\le CkF(k)\qquad(k\ge k_0),
\end{equation}
where $C\ge1$ is fixed and $k_0$ is sufficiently large.

\begin{lemma}\label{lem:words}
Put $a_0=(24C)^{-1}$. For every sufficiently large integer $\ell$
there are integers $r,s,t,v,a$ such that
\begin{equation}\label{eq:repeat}
 t=r+s\le\ell,\quad 0\le r<t,\quad s\ge v\ge a_0\ell/F(\ell),
 \quad 0\le a\le b^t-b^r,
\end{equation}
and
\begin{equation}\label{eq:approx}
 |(b^t-b^r)\xi-a|\le b^{-v}.
\end{equation}
\end{lemma}
\begin{proof}
Take $k=\lfloor\ell/(4CF(\ell))\rfloor$. For large $\ell$,
$k\ge k_0$, $k\ge\ell/(8CF(\ell))$, and $p(k)\le\ell/4$.
Among the $\ell-k+1$ length-$k$ blocks of the prefix there are two equal
ones, with starting-position difference $D$.
If $D\ge k$ they are disjoint. If $D<k$, the substring of length $k+D$
has period $D$, and its first two blocks of length
$v'=D\lfloor(k+D)/(2D)\rfloor$ are equal and disjoint.
Here $v'\ge k/3$: for $D\ge k/3$ use $v'\ge D$, and otherwise use
$v'>(k-D)/2$. Thus in either case the prefix has the form $UVWVX$
with $v=|V|\ge a_0\ell/F(\ell)$.

Set $r=|U|$, $s=|VW|$, $t=r+s$,
$A_j=\sum_{i=1}^jd_ib^{j-i}$, $A_0=0$, and $a=A_t-A_r$.
The number $a/(b^t-b^r)$ has base-$b$ expansion $U(VW)^\infty$ and agrees
with $\xi$ through the first $t+v$ digits. Hence its distance from
$\xi$ is at most $b^{-t-v}$, which gives \eqref{eq:approx}.
Writing $P$ for the value of $VW$ gives
$a=(b^s-1)A_r+P$, so $0\le a\le b^t-b^r$.
The estimate remains valid when the periodic tail consists entirely
of the digit $b-1$, by using the closed interval of values with a
specified finite prefix.
The length conditions follow from disjointness.
\end{proof}

Fix one such choice for each $\ell$.
\begin{lemma}\label{lem:selection}
There are increasing integers $\ell_j$ and associated approximations
$r_j,s_j,t_j,v_j,a_j$ satisfying
\begin{equation}\label{eq:select}
 s_{j+1}>2t_j,\qquad t_{j+1}>2t_j,
\end{equation}
and
\begin{equation}\label{eq:growth}
 j\ll\log\ell_j\ll j\log(2j),\qquad
 \log\log\ell_j\asymp\log(2j).
\end{equation}
\end{lemma}
\begin{proof}
Choose $\ell_1$ large. Define $\ell_{j+1}$ as the first $\ell>\ell_j$
with $s_\ell>2t_j$. It exists since
$s_\ell\ge a_0\ell/F(\ell)\to\infty$.
For fixed $D>4/a_0$, \eqref{eq:slow} shows that
$\ell^*=\lceil D\ell_jF(\ell_j)\rceil$ satisfies
$s_{\ell^*}\ge(a_0D/2)\ell_j>2t_j$ for large $\ell_j$.
Thus $\ell_{j+1}\ll\ell_jF(\ell_j)$.
With $y_j=\log\ell_j$, this yields
$y_{j+1}\le y_j+u\log y_j+O(1)$, hence
$y_j\le A j\log(2j)$ by induction for a sufficiently large fixed $A$.
For example choose $A$ so that
$u\log(Aj\log(2j))+O(1)\le A\log(2j)$ for all $j\ge1$.
The reverse bound follows from $\ell_j\ge t_j>2^{j-1}t_1$.
Taking logarithms gives the last assertion of \eqref{eq:growth}.
\end{proof}

\section{Twisted heights and the exact weights}\label{sec:height}
For a number field $K$ write $M_K$ for its places, normalized by
$\|z\|_v=|z|_p^{d_v}$ on $z\in\Q$ when $v\mid p$, where
$d_v=[K_v:\Q_p]/[K:\Q]$ and $\sum_{v\mid p}d_v=1$.
At an infinite place represented by an embedding $\sigma_v$ the
corresponding formula is $\|z\|_v=|\sigma_v(z)|^{d_v}$.
For a system of linearly independent forms
$\LL=(L_{iv}:v\in M_K,1\le i\le3)$ and real weights $c_{iv}$, put
\begin{equation}\label{eq:height}
 H_{\LL,c,Q}(x)=\prod_{v\in M_K}
 \max_{1\le i\le3}\bigl(\|L_{iv}(x)\|_vQ^{-c_{iv}}\bigr).
\end{equation}
For $x\in\Qbar^3$ use a finite extension $E/K$ containing its coordinates,
keep the forms unchanged at $w\mid v$, and multiply the weights by
$d(w\mid v)=[E_w:K_v]/[E:K]$. This gives an extension-independent height.

Choose a finite Galois extension $K/\Q$ in $\mathbb C$ containing
$\xi$, and let $G=\operatorname{Gal}(K/\Q)$ and
$d=[\Q(\xi):\Q]\ge2$. Write $\beta_1,\ldots,\beta_d$ for the distinct
conjugates of $\xi$. They are all nonzero.
For each infinite place $v$, choose $\sigma_v\in G$ with
$\|z\|_v=|\sigma_v(z)|^{d_v}$, and set $\beta_v=\sigma_v^{-1}(\xi)$.
Define the fixed forms by
\begin{equation}\label{eq:forms}
 (L_{1v},L_{2v},L_{3v})=
 \begin{cases}(X_1,X_2,L_{\beta_v}),&v\mid\infty,\\
 (X_1,X_2,X_3),&v\nmid\infty,
 \end{cases}
 \quad L_\beta=-\beta X_1+\beta X_2+X_3.
\end{equation}
Their local determinants are all $1$, and there are $d+3$ distinct forms.
For every rational vector $x$,
\begin{equation}\label{eq:transport}
 \|L_{\beta_v}(x)\|_v
 =|\sigma_v(L_{\beta_v}(x))|^{d_v}=|L_\xi(x)|^{d_v}.
\end{equation}
This is the inverse-conjugate transport used in
\cite[Section~5, equations (5.14)--(5.17)]{BE08}.
It uses smallness only at the specified real value of $\xi$.

\begin{lemma}[Mass of a conjugate]\label{lem:mass}
For every conjugate $\beta$ of $\xi$,
\begin{equation}\label{eq:mass}
 \sum_{\substack{v\mid\infty\\\beta_v=\beta}}d_v=\frac1d.
\end{equation}
\end{lemma}
\begin{proof}
Let $\iota\in G$ be complex conjugation on $K$ and
$H=\langle\iota\rangle$. Infinite places are represented by the left
cosets $H\sigma$. Since $\xi$ is real,
$(\iota\sigma)^{-1}(\xi)=\sigma^{-1}(\xi)$, so $\beta_v$ does not depend
on the representative. The field $K$ is either totally real or totally
imaginary, and $d_v=|H|/|G|$ in either case.
The set of $\sigma\in G$ with $\sigma^{-1}(\xi)=\beta$ has cardinality
$|G|/d$ by orbit--stabilizer. It is a union of left $H$-cosets, each of
cardinality $|H|$. Its corresponding total place weight is therefore
$(|G|/(d|H|))(|H|/|G|)=1/d$.
\end{proof}

For each rational prime $p\mid b$, put
\begin{equation}\label{eq:primeweights}
 \omega_p=\frac{\operatorname{ord}_p(b)\log p}{\log b},\qquad
 \kappa_v=d_v\omega_p\quad(v\mid p).
\end{equation}
These numbers are positive and satisfy
$\sum_{p\mid b}\omega_p=1$ and
$\sum_{v\mid b}\kappa_v=1$, where $v\mid b$ means $v\mid p$ for some
$p\mid b$. Let $S$ consist of all infinite places and all places above
the prime divisors of $b$.

Fix $0<\eps\le1/2$ and $0\le\rho\le1-\eps$. In the order of
\eqref{eq:forms}, define raw weights by
\begin{equation}\label{eq:raw}
 e_v=\begin{cases}
 d_v(1,\rho,-\eps),&v\mid\infty,\\
 \kappa_v(-1,-\rho,0),&v\mid b,\\
 (0,0,0),&v\notin S.
 \end{cases}
\end{equation}
Put
\begin{equation}\label{eq:center}
 \overline e_v=\tfrac13\sum_i e_{iv},\quad D_\eps=1+\eps/3,
 \quad c_{iv}(\rho)=\frac{e_{iv}-\overline e_v}{D_\eps},
 \quad \delta=\frac{\eps}{3+\eps}.
\end{equation}
The two mass identities give
\begin{equation}\label{eq:norm}
 \sum_i c_{iv}=0,\qquad \sum_v\max_i c_{iv}=1,
 \qquad \sum_v\overline e_v=-\eps/3.
\end{equation}
Indeed $\sum_{v,i}e_{iv}=-\eps$ and the sum of the raw local maxima
is $1$; centering changes the latter to $1+\eps/3$.

For a fixed integer $J$, take $\eps=a_0/F(\ell_J)$ and $j\le J$.
By monotonicity of $F$ and \eqref{eq:repeat},
$v_j\ge\eps t_j$ and $t_j-r_j\ge v_j$, so
$\rho_j=r_j/t_j\le1-\eps$.
Set $x_j=(b^{t_j},b^{r_j},a_j)$, $\Psi_j=b^{t_j}$ and
$Q_j=\Psi_j^{D_\eps}$.
At the infinite places, \eqref{eq:approx} and \eqref{eq:transport} give
$\|L_{iv}(x_j)\|_v\le\Psi_j^{e_{iv}(\rho_j)}$.
At a place $v\mid p\mid b$, the same bounds follow from
\[
 \|b^{t_j}\|_v=\Psi_j^{-\kappa_v},\qquad
 \|b^{r_j}\|_v=\Psi_j^{-\kappa_v\rho_j},\qquad
 \|a_j\|_v\le1.
\]
At all other finite places the coordinates are integral. Hence
\begin{equation}\label{eq:small}
 H_{\LL,c(\rho_j),Q_j}(x_j)
 \le\Psi_j^{\sum_v\overline e_v}=Q_j^{-\delta},
 \qquad Q_{j+1}>Q_j^2.
\end{equation}
No primitivity of $x_j$ is needed; $\Psi_j$ is an external parameter,
not an identification with its projective height.

\section{Uniform semistability}\label{sec:stable}
For a $q$-dimensional subspace $U\subset\Qbar^3$, define its local
weight as the minimum of $\sum_{i\in I}c_{iv}$ over all $q$-element
sets $I$ for which $(L_{iv}|_U)_{i\in I}$ is linearly independent.
Its global weight $w_c(U)$ is the sum over places. The system is
semistable if $w_c(U)\le0$ for every nonzero proper $U$.
Define $w_e$ analogously. Centering gives
\begin{equation}\label{eq:weightshift}
 w_c(U)=\frac{w_e(U)+q\eps/3}{D_\eps}.
\end{equation}

\begin{proposition}\label{prop:stable}
For the family \eqref{eq:raw}--\eqref{eq:center}, all nonzero proper
subspaces have negative weight. More precisely,
\[
 w_c(U)\le-\frac{\eps}{6D_\eps}\quad(\dim U=1),\qquad
 w_c(U)\le-\frac{\eps}{3D_\eps}\quad(\dim U=2).
\]
\end{proposition}
\begin{proof}
By Lemma~\ref{lem:mass}, the infinite-place contribution to $w_e(U)$
can be computed using each distinct $L_\beta$ once, with raw weights
$(1,\rho,-\eps)/d$. The finite contributions combine into one minimum
calculation with coordinate forms and weights $(-1,-\rho,0)$, since
$\sum_{v\mid b}\kappa_v=1$. This grouping concerns only the subspace
weight calculation; all original places and norms remain in the height.

Consider a line $U$. If no $L_\beta$ vanishes on it, its infinite
contribution is $-\eps$ and its finite contribution is at most zero.
If two distinct $L_\beta$ vanish, subtracting them shows
$U=\langle(1,1,0)\rangle$; all $L_\beta$ then vanish and
$w_e(U)=\rho-1\le-\eps$.
Otherwise precisely one $L_\beta$ vanishes. If $X_1|_U\ne0$, the
exceptional infinite contribution is at most $1/d$, the remaining
infinite contribution is $-(1-1/d)\eps$, and the finite contribution
is $-1$. Thus $w_e(U)\le-(1-1/d)(1+\eps)$.
If $X_1|_U=0$, then $X_2|_U\ne0$ and
$w_e(U)=-(1-1/d)(\rho+\eps)$.
Since $d\ge2$, every line therefore has $w_e(U)\le-\eps/2$.

For planes write $U=\ker M$, and put
$A_\beta=\spanop(L_\beta,X_2)$ and $B=\spanop(X_1,X_2)$ in the dual
space. For distinct conjugates $\beta,\beta'$,
\[
 A_\beta\cap A_{\beta'}=A_\beta\cap B=\langle X_2\rangle.
\]
This uses $\beta\ne\beta'$ and $\beta\ne0$.
At a grouped infinite place the minimum pair has weight
$(\rho-\eps)/d$ unless $M\in A_\beta$; the finite minimum is
$-1-\rho$ unless $M\in B$.
The following list is exhaustive. Proportional normals are identified.
\begin{center}\small
\begin{tabular}{@{}ll@{}}\toprule
Position of $M$ & $w_e(\ker M)$\\\midrule
$M\notin B\cup\bigcup_\beta A_\beta$ & $-1-\eps$\\
$M\in A_\beta\setminus(\langle L_\beta\rangle\cup\langle X_2\rangle)$
 & $-1-\eps+(1-\rho)/d$\\
$M\in\langle L_\beta\rangle$ for some $\beta$ & $-(1-1/d)(1+\eps)$\\
$M\in B\setminus(\langle X_1\rangle\cup\langle X_2\rangle)$ & $\rho-1-\eps$\\
$M\in\langle X_1\rangle\cup\langle X_2\rangle$ & $-\eps$\\\bottomrule
\end{tabular}
\end{center}
Each entry is at most $-\eps$, since $d\ge2$, $\eps\le1/2$ and
$\rho\le1-\eps$.
For example, at $M=L_\beta$ the pair for that conjugate must be
$(X_1,X_2)$, whereas at $M=X_1$ the finite pair is $(X_2,X_3)$.
These observations also cover the ties when $\rho=0$.
Equation \eqref{eq:weightshift} proves both bounds.
\end{proof}

\section{A finite-support moving-weight bound}\label{sec:moving}
We prove the arithmetic input separately. In this section $K$, $S$, and
$\LL$ are fixed; the dimension is three. Assume $S$ is finite and contains
all infinite places, $\LL$ is the coordinate system outside $S$, and the
determinant at every place is $1$. Write $s=|S|$.
Let $v_0\notin S$ be a fixed finite place.
Weights called admissible below satisfy
\begin{equation}\label{eq:admissible}
 c_{iv}=0\ (v\notin S),\qquad
 \sum_i c_{iv}=0,\qquad \sum_v\max_i c_{iv}\le1,
 \qquad w_c(U)\le0\ (0\ne U\subsetneq\Qbar^3).
\end{equation}

\begin{theorem}\label{thm:sparse}
There are constants depending only on the fixed data with the following
property. For $0<\delta\le1$ there exist $\Omega(\delta)>2$ and
$C_*(\delta)>e$ with
\begin{equation}\label{eq:cutofforder}
 \Omega(\delta)\ll\delta^{-5},\qquad
 \log\log C_*(\delta)\ll\delta^{-2}\log(2/\delta),
\end{equation}
such that any sequence $(c^{(h)},Q_h)$ of admissible weights and parameters
with
\[
 Q_1\ge C_*(\delta),\qquad Q_{h+1}>Q_h^{\Omega(\delta)},\qquad
 \{x\ne0:H_{\LL,c^{(h)},Q_h}(x)\le Q_h^{-\delta}\}\ne\varnothing
\]
has length $O(\delta^{-2})$.
\end{theorem}

We first isolate the pointwise part of the proof in \cite{EF13}. This is
needed to control the cutoff, rather than import the larger constant
in their final interval theorem.

\begin{lemma}[Pointwise exterior data]\label{lem:packet}
There are constants $A_0>0$ and $c_0>0$, depending only on $K,\LL$, such
that if $c$ is admissible, $Q\ge\exp(A_0/\delta)$ and
$H_{\LL,c,Q}(x)\le Q^{-\delta}$ for some nonzero $x$, then there is
$k\in\{1,2\}$ with the following properties. Put $N=\binom3k=3$.
There are fixed exterior forms $\wt L_{lv}$, depending only on $k,\LL$,
and independent vectors $z_1,z_2\in\Qbar^3$ spanning a plane $T$ such that
\begin{equation}\label{eq:packetheight}
 H_2(T)\ge Q^{c_0\delta}.
\end{equation}
For every place $w$ of a common finite extension $E/K$,
\begin{equation}\label{eq:packetlocal}
 \|\wt L_{lw}(z_j)\|_w\le
 \begin{cases}Q^{d(w\mid v)a_{lv}},&w\mid v\ne v_0,\\
 Q^{d(w\mid v_0)b_l},&w\mid v_0,
 \end{cases}
\end{equation}
where, with $\gamma_v=\max_i c_{iv}$,
\begin{equation}\label{eq:packetweights}
 \sum_l a_{lv}=0,\quad |a_{lv}|\le2\gamma_v,\quad
 a_{lv}=0\ (v\notin S),\quad
 \sum_l b_l\le-\delta/3,\quad |b_l|\le3.
\end{equation}
Here $H_2$ is the absolute Euclidean height of the exterior coordinate
vector defining a subspace, as in \cite[Section~6]{EF13}.
\end{lemma}
\begin{proof}
Let $\lambda_1\le\lambda_2\le\lambda_3$ be the successive infima of the
twisted height over $\Qbar$. The fixed-form comparison and the absolute
Minkowski theorem \cite[Lemma~7.1 and Proposition~9.2]{EF13} give
\begin{equation}\label{eq:minima}
 C_0^{-1}Q^{-1}\le\lambda_1,\qquad
 3^{-3/2}\le\lambda_1\lambda_2\lambda_3\le8.
\end{equation}
The product bounds use the determinant assumption. As
$\lambda_1\le Q^{-\delta}$, enlarging $A_0$ ensures $\lambda_3\ge1$.
One of the two consecutive ratios is therefore at most $Q^{-\delta/2}$;
fix such a $k$. Its infimum subspace $T_k$ has dimension $k$.

For clarity we also give the height estimate behind this assertion.
Choose a basis $g_1,\ldots,g_k$ with twisted heights arbitrarily close
to $\lambda_1,\ldots,\lambda_k$, and at each place select an independent
restricted $k$-tuple realizing its local minimum weight.
Semistability and the determinant estimate imply that the product of
these local determinant norms is at most
$C\lambda_1\cdots\lambda_k\le C'Q^{-\delta/3}$.
The last inequality follows from \eqref{eq:minima} and the gap, since
$k(3-k)/6=1/3$ for $k=1,2$.
Let $R\ge3$ bound the number of distinct original forms. Their nonzero
$k$-fold determinants on the basis number at most $\binom Rk$.
The product formula, applied to these determinants, gives the lower
bound $(C''H_2(T_k))^{-(\binom Rk-1)}$ for the same product; this is
also \cite[Lemma~10.2]{EF13}. The fixed constants can be absorbed when
$\log Q\ge A_0/\delta$, giving
$H_2(T_k)\ge Q^{c_0\delta}$, for example with
$c_0=1/(12R^3)$ after increasing $A_0$.

Apply Davenport's construction \cite[Lemma~11.3]{EF13} separately to
this height and take the $(3-k)$th exterior power as in its
equations (11.10)--(11.17). The exterior forms depend only on $k,\LL$.
Away from $v_0$ the exponents $a_{lv}$ are sums of $3-k$ original weights;
their sums are zero and $|a_{lv}|\le2\gamma_v$.
The first $N-1=2$ exterior vectors span $T$, whose height equals
$H_2(T_k)$ by \cite[Lemma~6.1]{EF13}.

Here is the size check at the auxiliary place, including the cutoff.
Let $\nu_l$ be the products of $3-k$ distinct successive infima.
Up to a permutation, the bounds defining $Q^{b_l}$ are
$3^{27}\nu_1,\ldots,3^{27}\nu_{N-1},3^{27}\nu_{N-1}$.
Thus their product is
\[
 3^{27N}(\lambda_1\lambda_2\lambda_3)^{\binom2{2-k}}
       \frac{\lambda_k}{\lambda_{k+1}}
 \le C Q^{-\delta/2}\le Q^{-\delta/3}.
\]
Also \eqref{eq:minima} bounds any product of $p=1,2$ distinct infima
between $C^{-1}Q^{-p}$ and $C Q^{3-p}$. It follows that $|b_l|\le3$
once $A_0$ is sufficiently large. All absorptions above have the form
$Q^{a\delta}\ge C$ with fixed $a,C>0$, or $Q^a\ge C$.
Consequently $\exp(A_0/\delta)$ suffices throughout; the final cutoff
of \cite[Theorem~8.1]{EF13} has not been used.
\end{proof}

\subsection{The auxiliary polynomial with changing weights}
Suppose $m$ of the parameters have the same $k$ in
Lemma~\ref{lem:packet}; label them $h=1,\ldots,m$.
Their exterior forms are consequently the same.
Write $a_{hlv},b_{hl}$ for their exterior weights and set
\begin{equation}\label{eq:Gamma}
 \Gamma_v=\max_{1\le h\le m}\max_i c^{(h)}_{iv}.
\end{equation}
Each $\Gamma_v\ge0$, and admissibility gives
\begin{equation}\label{eq:Gammasum}
 \sum_{v\in S}\Gamma_v\le s.
\end{equation}
All ordinary-place vanishing conditions below are imposed only at
places with $\Gamma_v>0$.

Let $r_1,\ldots,r_m$ be positive integers, $D_r=\sum_h r_h$, and let
$\mathcal U(r)$ consist of arrays $j=(j_{hl})$ of nonnegative integers
with $\sum_l j_{hl}=r_h$ for each $h$. Put $V=|\mathcal U(r)|$.
For $0<\alpha<1$ satisfying
\begin{equation}\label{eq:mcondition}
 m>2\alpha^{-2}\log(2s+2),
\end{equation}
there is a nonzero multihomogeneous polynomial $P$ over $K$, of block
degrees $r_h$, with the following properties. Expand $P$ in the exterior
forms at a place $v$, writing its coefficients as $d^{(v)}_j$.
Then
\begin{align}
 d^{(v)}_j&=0 &&\text{if }\sum_{h,l}j_{hl}a_{hlv}/r_h\ge3m\alpha\Gamma_v,
       \quad v\in S,\ \Gamma_v>0,\label{eq:vanishordinary}\\
 d^{(v_0)}_j&=0 &&\text{if }\sum_{h,l}j_{hl}b_{hl}/r_h\ge-m\delta/9+3m\alpha,
       \label{eq:vanishspecial}\\
 H_2(P)&\le C_K A^{D_r}.&&\label{eq:polyheight}
\end{align}
Here and below $A\ge2$ and $C_K\ge1$ are fixed constants independent of
$m,r,c^{(h)},Q_h$.

To prove this, choose $j$ uniformly from $\mathcal U(r)$. The blocks are
independent, and
$r_h^{-1}\sum_l j_{hl}a_{hlv}$ has mean zero and lies in
$[-3\Gamma_v,3\Gamma_v]$.
Hoeffding's inequality therefore bounds the fraction satisfying the
condition in \eqref{eq:vanishordinary} by $e^{-m\alpha^2/2}$.
For the special place the mean of the total is at most $-m\delta/9$
and the individual summands lie in $[-3,3]$, giving the same bound.
This is precisely the block-dependent form of
\cite[Lemma~13.3]{EF13}.
The total number of imposed linear equations in the $V$ coefficients
of $P$ is less than $V/2$ by \eqref{eq:mcondition}.
The changes of coordinates use fixed matrices. Consequently each row
of this linear system has height at most $A_1^{D_r}$ and
$V\le(eN)^{D_r}$, uniformly in which equations were selected.
The number-field Siegel lemma \cite[Lemma~13.1]{EF13}, with exponent
$U/(V-U)<1$, gives \eqref{eq:polyheight} after increasing $A$.
This proves all three assertions, without discretizing places or weights.

For a derivative multi-index $i$, use divided derivatives
$P_i=(\prod_{h,l}i_{hl}!)^{-1}\partial^iP$ and write
\[
 |i|_r=\sum_{h,l}i_{hl}/r_h.
\]
If $|i|_r\le2m\alpha$, expansion in the exterior forms gives
\begin{align}
 P_i&=\sum_jd^{(v)}_{i,j}\prod_{h,l}\wt L_{lv}(X_h)^{j_{hl}},
       \label{eq:derivativeexp}\\
 d^{(v)}_{i,j}&=0\quad\text{if }\sum_{h,l}j_{hl}a_{hlv}/r_h>12m\alpha\Gamma_v,
       \quad v\in S,\label{eq:derivativeordinary}\\
 d^{(v_0)}_{i,j}&=0\quad\text{if }\sum_{h,l}j_{hl}b_{hl}/r_h>
                  -m\delta/9+12m\alpha,\label{eq:derivativespecial}\\
 \prod_v\max_j\|d^{(v)}_{i,j}\|_v&\le C_K A_2^{D_r}.
       \label{eq:derivativeheight}
\end{align}
Indeed a contributing original monomial has index $j+q$, where
$\sum_lq_{hl}=\sum_li_{hl}$ in every block. Its weighted shift has
absolute value at most $6m\alpha\Gamma_v$, or $6m\alpha$ at $v_0$.
Either strict inequality above would therefore force that original
coefficient to vanish by \eqref{eq:vanishordinary} or
\eqref{eq:vanishspecial}. The coefficient bound follows from fixed
inverse coordinate matrices and binomial coefficient bounds, exactly
as in \cite[Lemma~13.4]{EF13}; combined with \eqref{eq:polyheight} these
give a fixed $A_2$. At places outside $S\cup\{v_0\}$ all the weights
are zero, so no vanishing condition is required.

\subsection{Nonvanishing, the product formula, and the cutoff}
We now prove Theorem~\ref{thm:sparse}. Choose
\begin{equation}\label{eq:parameters}
 \alpha=\frac{\delta}{1000(s+1)},\qquad
 m=\left\lfloor2\alpha^{-2}\log(2s+2)\right\rfloor+1,
 \qquad \Omega=\frac{4m^2}{\alpha}.
\end{equation}
There are only two possible $k$, so a sequence of at least $2m-1$
parameters has a subsequence of $m$ parameters with the same $k$.
Separation is preserved. Suppose such a subsequence exists and label
it by $Q_1<\cdots<Q_m$.

Choose an arbitrarily large positive integer $r_1$ with
$r_1\log Q_1>\alpha^{-1}\log Q_m$ and $C_K\le2^{r_1}$, and put
\[
 r_h=\left\lceil\frac{r_1\log Q_1}{\log Q_h}\right\rceil\ (h>1),
 \qquad B=Q_1^{r_1}.
\]
Then
\begin{equation}\label{eq:degrees}
 B\le Q_h^{r_h}<B^{1+\alpha},\quad r_h\le r_1,\quad
 \frac{r_h}{r_{h+1}}>\frac{\Omega}{1+\alpha}\ge\frac{2m^2}{\alpha}.
\end{equation}
In particular the decreasing degree ratio has the direction required
by \cite[Proposition~12.1, equation (12.3)]{EF13}.
Construct $P$ as above. For each exterior plane $T_h$, use the grid
\[
 \mathcal G_h=\{u_1z_{h1}+u_2z_{h2}:u_j\in\Z,
                    |u_j|\le\lceil N/\alpha\rceil\},\qquad N=3.
\]
Put $E_m=(N-1)(3m^2/\alpha)^m$ and $A_3=2eA$.
By \eqref{eq:polyheight} and $C_K\le2^{r_1}$,
$e^{D_r}H_2(P)\le A_3^{D_r}$.
The other hypothesis of that nonvanishing proposition is
\begin{equation}\label{eq:nonvanishheight}
 H_2(T_h)^{r_h}\ge(e^{D_r}H_2(P))^{E_m}.
\end{equation}
By \eqref{eq:packetheight} and \eqref{eq:degrees} it is enough that
\begin{equation}\label{eq:cutoffone}
 c_0\delta\log Q_1\ge mE_m\log A_3.
\end{equation}
The proposition then supplies $x_h\in\mathcal G_h\setminus\{0\}$ and
$|i|_r\le2m\alpha$ such that $P_i(x_1,\ldots,x_m)\ne0$.

We estimate this value. For an infinite place $w$, let $s(w)$ be its
local degree divided by $[E:\Q]$, and set $s(w)=0$ at finite places.
The grid and \eqref{eq:packetlocal} give the extra factor
\begin{equation}\label{eq:gridfactor}
 \|\wt L_{lw}(x_h)\|_w\le
 (N^2/\alpha)^{s(w)}Q_h^{d(w\mid v)a_{hlv}}\quad(w\nmid v_0),
\end{equation}
with the same bound at $w\mid v_0$ using $b_{hl}$ and no grid factor.
The factor $N^2/\alpha$ is retained throughout.

For a surviving monomial at $w\mid v\ne v_0$, put
$A_{hw}=r_h^{-1}\sum_lj_{hl}d(w\mid v)a_{hlv}$ and
$\Gamma_w=d(w\mid v)\Gamma_v$.
The derivative degrees are at most $r_h$, so $|A_{hw}|\le3\Gamma_w$.
Replacing $Q_h^{r_h}$ by $B$ costs at most $3m\alpha\Gamma_w$ in its
exponent, by \eqref{eq:degrees}.
Equation \eqref{eq:derivativeordinary} bounds the remaining sum by
$12m\alpha\Gamma_w$. The resulting exponent is at most
$15m\alpha\Gamma_w$.
At $w\mid v_0$ the same argument using \eqref{eq:derivativespecial}
gives at most $d(w\mid v_0)(-m\delta/9+15m\alpha)$.

The number of monomials in a derivative is at most
$V\le(eN)^{D_r}$. Multiplying the local bounds, using
\eqref{eq:derivativeheight}, $\sum_ws(w)=1$,
$\sum_{w\mid v}d(w\mid v)=1$, and \eqref{eq:Gammasum}, gives
\begin{align}
 \prod_w\|P_i(x_1,\ldots,x_m)\|_w
 &\le G(\alpha)^{D_r}
       (B^m)^{-\delta/9+15(s+1)\alpha}\notag\\
 &\le G(\alpha)^{D_r}(B^m)^{-\delta/12},\label{eq:product}
\end{align}
where $G(\alpha)=2A_2eN^3/\alpha$. We used
$C_K\le2^{r_1}\le2^{D_r}$.
Since $D_r\le mr_1$, the right side is strictly less than $1$ if
\begin{equation}\label{eq:cutofftwo}
 \log Q_1\ge24\delta^{-1}\log G(\alpha).
\end{equation}
This contradicts the product formula for the nonzero value of $P_i$.

For definiteness, all the required conditions hold with
\begin{equation}\label{eq:explicitcutoff}
 C_*(\delta)=\exp\left(
 \frac{A_0+1}{\delta}
 +\frac{2mE_m\log A_3}{c_0\delta}
 +\frac{24\log G(\alpha)}{\delta}\right).
\end{equation}
As $m\ll\delta^{-2}$ and $\alpha\asymp\delta$,
$\Omega\ll\delta^{-5}$ and
\[
 \log\log C_*(\delta)
 \ll m\log(m/\alpha)+\log(2/\delta)
 \ll\delta^{-2}\log(2/\delta).
\]
The maximum length is at most $2m-2$. This proves the theorem.

\begin{corollary}\label{cor:count}
Under the same hypotheses, suppose $Q_{h+1}>Q_h^2$.
The number of parameters at or above $C_*(\delta)$ is
$O(\delta^{-2}\log(2/\delta))$.
\end{corollary}
\begin{proof}
Set $q=\lceil\log_2\Omega\rceil\ll\log(2/\delta)$ and partition the
indices modulo $q$. In each part consecutive parameters satisfy the
separation condition of Theorem~\ref{thm:sparse}. Apply its bound to
each part.
\end{proof}

\section{Completion of the complexity argument}\label{sec:completion}
Apply the moving-weight theorem with the fixed system \eqref{eq:forms}
and the finite set $S$ from Section~\ref{sec:height}.
Choose a fixed rational prime $p_0\nmid b$ and a place above it as $v_0$.
All admissibility conditions follow from \eqref{eq:norm} and
Proposition~\ref{prop:stable}.
For a prefix sequence of length $J$ constructed under \eqref{eq:LC},
use the common $\eps=a_0/F(\ell_J)$ and
$\delta=\eps/(3+\eps)$. These depend on $J$, whereas the forms, field,
and support do not. Hence every implied constant in the arithmetic
bound is uniform in $J$, although it may depend on $\xi$ and $b$.

Let $y_J=\log\ell_J$. By \eqref{eq:growth},
\begin{equation}\label{eq:budget}
 \delta^{-2}\log(2/\delta)
 \ll F(\ell_J)^2\log(2F(\ell_J))
 \ll J^{2u}(\log(2J))^{2u+1-2\gamma}.
\end{equation}
For the last step use both $J\ll y_J\ll J\log(2J)$ and
$\log y_J\asymp\log(2J)$, so that negative powers of $\log y_J$ are
also controlled.
If $u<1/2$ and $\gamma=0$, the last expression is $o(J)$.
If $u=1/2$ and $\gamma>1$, it is
$J(\log(2J))^{2-2\gamma}=o(J)$ as well.

In either case \eqref{eq:cutofforder} implies
$\log\log C_*(\delta)=o(J)$.
But \eqref{eq:select} and $Q_j=b^{D_\eps t_j}$ imply
\[
 \log\log Q_j=\log t_j+O(1)\ge(j-1)\log2+O(1).
\]
Consequently every $j$ with $\lceil J/2\rceil\le j\le J$ has
$Q_j\ge C_*(\delta)$ for sufficiently large $J$.
Corollary~\ref{cor:count}, together with \eqref{eq:small}, now gives
\[
 J/2\ll\delta^{-2}\log(2/\delta)=o(J),
\]
a contradiction. Taking $u=1/2$ proves \eqref{eq:main}, since a finite
limsup would imply \eqref{eq:LC} for some $C$.
For $0<\eta<1/2$ the direct argument with $u=\eta$, $\gamma=0$ proves
\eqref{eq:eta}. Alternatively \eqref{eq:main} implies it for every
real $\eta<1/2$, because
$(\log L)^{1/2-\eta}/(\log\log L)^\gamma\to\infty$.
This completes the proof of Theorem~\ref{thm:main}.

\section{Scope of the argument}
The improvement uses two distinct facts: the finite common support
permits simultaneous coefficient vanishing with $m=O(\delta^{-2})$,
and the entire family of exact weights is semistable. The first does
not follow just by applying a fixed-weight theorem at each point.
The second is a calculation about every algebraic subspace, rather
than only the rational approximating vectors.

The proof deliberately uses the cutoff in \eqref{eq:explicitcutoff}.
Using the cutoff from the published final interval theorem without
rechecking its preliminary lemmas would introduce an avoidable
additional logarithm in the double-logarithmic threshold estimate.
Similarly, the grid factor in \eqref{eq:gridfactor} is harmless only
after its contribution is included in that cutoff.

Lemma~\ref{lem:mass} and the normalization
$\sum_{v\mid b}\kappa_v=1$ make the semistability argument applicable
to every real algebraic irrational number and every integer base $b\ge2$.
Nonreal conjugates and splitting or ramification of primes cause no
additional loss in the exponent. The finite support, and consequently
the constants, are allowed to depend on the fixed pair $(\xi,b)$.
No assertion uniform in a varying base or varying algebraic number is
made about these constants.

\begin{thebibliography}{9}
\bibitem{AB07} B.~Adamczewski and Y.~Bugeaud,
On the complexity of algebraic numbers I. Expansions in integer bases,
\emph{Ann. of Math.} (2) \textbf{165} (2007), no.~2, 547--565.
\href{https://doi.org/10.4007/annals.2007.165.547}{doi:10.4007/annals.2007.165.547}.
\bibitem{BE08} Y.~Bugeaud and J.-H.~Evertse,
On two notions of complexity of algebraic numbers,
\emph{Acta Arith.} \textbf{133} (2008), no.~3, 221--250.
\href{https://doi.org/10.4064/aa133-3-3}{doi:10.4064/aa133-3-3}.
\bibitem{EF13} J.-H.~Evertse and R.~G.~Ferretti,
A further improvement of the quantitative Subspace Theorem,
\emph{Ann. of Math.} (2) \textbf{177} (2013), no.~2, 513--590.
\href{https://doi.org/10.4007/annals.2013.177.2.4}{doi:10.4007/annals.2013.177.2.4}.
\end{thebibliography}
\end{document}
